\documentclass[UTF8,11pt]{ctexart}

\usepackage[a4paper,margin=2.4cm]{geometry}
\usepackage{amsmath,amssymb,bm}
\usepackage{physics}
\usepackage{hyperref}
\usepackage{booktabs}

\hypersetup{
    colorlinks=true,
    linkcolor=blue,
    citecolor=blue,
    urlcolor=blue
}

\title{La$_3$Ni$_2$O$_7$ RKKY }
\author{}
\date{\today}

\begin{document}

\maketitle

\section{RKKY}

非相互作用哈密顿量:
\begin{equation}
    H_0
    =
    \sum_{\mathbf k,\sigma}
    c_{\mathbf k a\sigma}^{\dagger}
    h_{ab}(\mathbf k)
    c_{\mathbf k b\sigma}.
\end{equation}


局域磁矩与巡游电子之间的交换耦合可以写成
\begin{equation}
    H_K
    =
    J_K
    \sum_i
    \mathbf S_i \cdot \mathbf s_i,
\end{equation}

在 $J_K$ 较小的情况下，对 $H_K$ 做二阶微扰，可以得到有效的 RKKY 相互作用
\begin{equation}
    \Delta E^{(2)}
    =
    \sum_{n\neq 0}
    \frac{
        \left|
        \langle n | H_K | 0 \rangle
        \right|^2
    }{
        E_0-E_n
    },
\end{equation}
\begin{equation}
    \Delta E^{(2)}
    =
    -
    J_K^2
    \sum_{i,j}
    \sum_{\alpha,\beta}
    S_i^\alpha S_j^\beta
    \sum_{n\neq 0}
    \frac{
        \langle 0 | s_i^\alpha | n \rangle
        \langle n | s_j^\beta | 0 \rangle
    }{
        E_n-E_0
    }.
\end{equation}
\begin{equation}
    \chi_{ij}^{\alpha\beta}
    =
    \sum_{n\neq 0}
    \frac{
        \langle 0 | s_i^\alpha | n \rangle
        \langle n | s_j^\beta | 0 \rangle
    }{
        E_n-E_0
    }.
\end{equation}
\begin{equation}
    H_{\mathrm{RKKY}}
    =
    \sum_{i,j}
    J_{ij}^{\mathrm{RKKY}}
    \,
    \mathbf S_i \cdot \mathbf S_j .
\end{equation}

\begin{equation}
    J_{ij}^{\mathrm{RKKY}}
    \propto
    -
    J_K^2
    \chi_{ij}^{s}.
\end{equation}


\begin{equation}
    \chi^{0}_{ab;cd}(\mathbf q)
    =
    -
    \frac{1}{N}
    \sum_{\mathbf k}
    \sum_{m,n}
    u_{c m}(\mathbf k+\mathbf q)
    u^{*}_{a m}(\mathbf k+\mathbf q)
    u_{b n}(\mathbf k)
    u^{*}_{d n}(\mathbf k)
    \frac{
        f[\varepsilon_m(\mathbf k+\mathbf q)-\mu]
        -
        f[\varepsilon_n(\mathbf k)-\mu]
    }{
        \varepsilon_m(\mathbf k+\mathbf q)
        -
        \varepsilon_n(\mathbf k)
    }.
    \label{eq:multi_orbital_chi}
\end{equation}
\begin{equation}
    \chi^{0}_{ab;cd}(\mathbf R)
    =
    \frac{1}{N_q}
    \sum_{\mathbf q}
    e^{i\mathbf q\cdot \mathbf R}
    \chi^{0}_{ab;cd}(\mathbf q).
\end{equation}
\end{document}
